\section{问题分析}

\subsection{问题一的分析}

针对问题一，首先，将一天划分为144个10\,min调度时段，并将负载功率和光伏预测功率转换为对应电量，以各时段的外部电网购电量和储能充放电量为主要决策变量。其次，根据各时段电量的来源与去向建立电量平衡关系，并依据充放电效率描述储电量在相邻时段间的变化，同时考虑储电量上下限、最大充放电功率、充放电互斥及日初日末储电量相同等约束。在此基础上，以全天购电费用最小为目标建立\textbf{混合整数线性规划模型}，最终求得满足题目各项运行要求的购电与储能调度方案。



\subsection{问题二的分析}

针对问题二，需在目标日实际负载和光伏出力未知的条件下提前制定购电计划，并考虑预测偏差可能产生的高价紧急购电。首先，利用附件2中的实际负载与实际光伏数据，将问题一的储能调度模型扩展至全年，求得完全信息条件下的最低购电费用，作为评价日前购电策略的理论基准；其次，利用目标日前可获得的历史数据预测负载与光伏出力，通过误差比较确定基础预测方法，并将预测结果代入混合整数线性规划模型，制定日前购电与储能调度方案；最后，进一步考虑星期和季节特征，结合近期净负荷偏差与日初储电量动态设置安全冗余，以降低预测偏差造成的紧急购电风险，形成最终的日前购电策略。

\subsection{问题三的分析}

问题三在问题二的基础上引入日内多时刻光伏预报，利用更新的预测信息调整尚未执行时段的购电量，并考虑购电调整费用与紧急购电费用。首先，沿用问题二的负载预测方法，结合 $0{:}00$ 获得的光伏发电功率预报建立混合整数线性规划模型，在满足供需平衡及储能运行约束的条件下，制定当天的计划购电与储能充放电方案；其次，将附件3提供的整点光伏预报转换至 $10\,\mathrm{min}$ 时间尺度，在 $6{:}00$、$12{:}00$ 和 $18{:}00$ 获得新预报后，保持已执行购电量和 $0{:}00$ 确定的储能充放电方案不变，对尚未执行时段的购电量进行调整，并比较调整前后的预计费用。


\subsection{问题四的分析}

问题四进一步考虑外部电网电价随日期和时段变化，需要在波动电价条件下重新制定问题二和问题三的购电策略。首先，利用目标日前可获得的历史电价预测当天各时段价格，并比较不同电价预测方法；其次，在问题二模型中引入预测电价，沿用负载与光伏预测方法及储能运行约束，重新求解混合整数线性规划模型，并根据实际电价计算紧急购电费用和总购电费用；再次，在问题三模型中同时引入光伏预报更新与电价动态更新，在 $6{:}00$、$12{:}00$ 和 $18{:}00$ 对尚未执行时段的购电量进行滚动调整，并比较不同电价预测方法和储能调整方式；最后，比较不同策略的紧急购电量与总购电费用，确定波动电价下的购电策略


\begin{figure}[H]
    \centering
    \includegraphics[width=1\linewidth]{图01_微网购电与储能调度的整体建模框架.png}
    \caption{微网购电与储能调度的整体建模框架}
    \label{fig:01}
\end{figure}